\begin{frame}\frametitle{Ответ рецензенту: DIR-24-8}
\footnotesize
\textbf{Комментарий:} рассмотреть DIR-24-8 и варианты для IPv6 как структуры,
используемые для FIB/RIB в обработчиках пакетов.

\vspace{0.3em}
\begin{itemize}\setlength{\itemsep}{0.35em}
  \item DIR-24-8 оптимизирован для dataplane-сценария: очень быстрый LPM
  по адресу при высокой доле операций чтения.
  \item В работе целевой профиль другой: микросервис поддерживает
  \texttt{prefix -> ASN}, начальный снимок, поток \texttt{Update}/\texttt{Withdraw}
  и полную выгрузку для внешних потребителей.
  \item Структура оценивалась не только по LPM, но и по стоимости точечных
  изменений, ресинхронизации и единой поддержке IPv4/IPv6.
  \item DIR-подход можно рассматривать как перспективное направление для
  read-optimized backend или кэша LPM-запросов.
\end{itemize}
\end{frame}

\begin{frame}\frametitle{Ответ рецензенту: ключ ART}
\footnotesize
\textbf{Комментарий:} почему ключ ART имеет вид
\texttt{семейство адресов, длина префикса, префикс}.

\vspace{0.3em}
\begin{itemize}\setlength{\itemsep}{0.35em}
  \item Такой ключ является каноническим идентификатором записи и однозначно
  различает IPv4, IPv6 и разные длины маски для одного адресного начала.
  \item \texttt{Upsert}/\texttt{Delete} становятся точечными операциями:
  сервис строит тот же ключ, выбирает шард и изменяет одну запись.
  \item LPM реализуется как набор точных поисков по активным длинам префикса;
  это сохраняет корректность и упрощает обработку BGP-событий.
  \item Ветвление на верхних уровнях является осознанным компромиссом:
  приоритет отдан изменяемости и потоковым обновлениям, а не максимальной
  оптимизации одиночного LPM.
\end{itemize}
\end{frame}

\begin{frame}\frametitle{Ответ рецензенту: шардирование}
\footnotesize
\textbf{Комментарий:} почему выбрана схема шардирования, если lookup может
обращаться к нескольким шардам.

\vspace{0.3em}
\begin{itemize}\setlength{\itemsep}{0.35em}
  \item Шардирование выбрано для локализации \texttt{Upsert}/\texttt{Delete}:
  обновление одного префикса блокирует только один ART, а не всё хранилище.
  \item Для BGP-потока это важнее, чем ускорение одиночного последовательного
  \texttt{Lookup}: события приходят независимо и должны применяться точечно.
  \item Эксперименты показывают, что в последовательных сценариях рост числа
  шардов не даёт выигрыша; поэтому число шардов является параметром настройки.
  \item Практическая схема: один шард для минимальных накладных расходов или
  умеренное шардирование при конкурентной нагрузке.
\end{itemize}
\end{frame}
